\section{Recovery of an uncalibrated homothety ratio}
\label{sec:unknown-scale}
Integrating a pooled forward collision mean eliminates the unknown
launch density and measures a sum of widths of the original obstacle.
Together with the two positivity supports, this scalar identifies the
ratio of two homothetic settings. An independent normalization uses
the area of a recovered occupation component and requires a weaker
separation condition.

Write $A$ for the physical footprint at the first setting, and let
$Z$ have an unknown probability density $j$ on $A$. The two settings
use the displacements $Z$ and $rZ$, where
\begin{equation}\label{eq:unk-scale-prior}
             1+g_*\le r\le R_*,\qquad g_*>0.
\end{equation}
Only these bounds, not the value of $r$, are supplied. The common
homothety origin is fixed in the nominal laboratory coordinates.
The footprint satisfies \eqref{eq:unknown-K-prior}, with $A$ in place
of $K$, and the density satisfies \eqref{eq:unknown-K-density}.
Assume $t+2R_*R_K<d_0$. Thus the two expanded-component collections
have a known positive separation greater than $t$, and their
cross-setting nesting margin is at least $g_*r_K$.

If the apparatus is originally described by unknown factors
$\lambda_1,\lambda_2$ acting on a parameter footprint $K$, then
$A=\lambda_1K$ and $r=\lambda_2/\lambda_1$. This is an operational
normalization: multiplying $K$ by a positive constant and dividing
both factors by that constant leaves the physical experiment unchanged.
The quantities recovered below are the physical first footprint and
the ratio, without assigning an additional unit to $K$.

All support functions in this section use laboratory coordinates.
For a unit vector $u$, $p_H(u)$ denotes $\sup_{x\in H}x\cdot u$,
the same support function expressed on the unit circle.

\subsection{An integrated collision width}
For a convex body $H$ write
\[
 \mathcal W(H)=p_H(e_1)+p_H(-e_1)+p_H(e_2)+p_H(-e_2),
\]
the sum of its two coordinate widths. This functional is invariant
under translation and additive under Minkowski addition. Let $A$ be
the first launch footprint, let $j$ be its probability density, and
put $P_{1,C}=C+(-A)$. Denote the first-setting pooled forward mean by
$F_1$.

\begin{lemma}[An isolated pooled flux measures width]
\label{lem:unk-scale-flux}
Suppose $2t+\diam A<d_0$. After sufficiently accurate fixed coarse
acquisition of a complete $P_{1,C}$, one can construct a bounded
domain $E$, consisting of finitely many squares of a fixed dyadic
mesh, which contains $P_{1,C}+t\overline B$ and no forward-response
support belonging to another component. Then
\begin{equation}\label{eq:unk-scale-flux}
 \Phi_C:=\int_E F_1(x)\dd x
                      =\frac t2\mathcal W(C).
\end{equation}
The domain and its area are computed from the coarse observations and
the fixed priors. Neither $j$ nor an occupation integral is an input.
\end{lemma}
\begin{proof}
For a displacement $a=tv$, $|v|=1$, set
$S_a(C)=C+[-a,0]$. A line parallel to $v$ meets $C$ in an interval,
a point, or the empty set. On every nonempty interior slice,
$S_a(C)\setminus C$ is an interval of length $t$. The transverse
projection of $C$ has length
$p_C(v^\perp)+p_C(-v^\perp)$. Cavalieri's formula therefore gives
\begin{equation}\label{eq:unk-scale-strip-area}
 \int_{\mathbb R^2}
       (\one_{S_a(C)}-\one_C)
   =t\{p_C(v^\perp)+p_C(-v^\perp)\}.
\end{equation}
Endpoint and tangency conventions affect only null sets. This is also
the segment case of the classical mixed-area formula \cite{Schneider}.

Since $t<d_0$, one segment of length $t$ cannot meet two distinct
physical components, and a start in one component cannot reach another.
Consequently the collision indicator decomposes pointwise as
\[
 B_a^{\mathcal O}
       =\sum_C(\one_{S_a(C)}-\one_C).
\]
The sum is locally finite and its summands are nonnegative. The
contribution $F_{1,C}$ obtained by averaging a single summand against
$j$ and the compass law is supported in
\[
 \bigcup_{v\in\{\pm e_1,\pm e_2\}}
       \{P_{1,C}+[-tv,0]\}
                    \subset P_{1,C}+t\overline B.
\]
Distinct expanded components are separated by at least
$d_0-\diam A$: offsets from the same footprint differ by at most its
diameter. Their displayed enlarged envelopes are therefore separated
by at least $d_0-\diam A-2t>0$. A sufficiently accurate coarse hull,
a fixed enlargement, and a sufficiently small fixed dyadic square mesh
give a domain $E$ containing the complete envelope of $C$ while
remaining disjoint from every other envelope. All error reserves use
fixed fractions of this positive separation. Thus $F_1=F_{1,C}$ on
$E$, and $F_{1,C}=0$ outside $E$.

Tonelli's theorem, translation invariance of area, and $\int_Aj=1$
now show that $\int_EF_1$ is the average of the four strip areas in
\eqref{eq:unk-scale-strip-area}. Each coordinate width occurs twice,
which proves \eqref{eq:unk-scale-flux}. No density differentiability,
upper bound, or symmetry is used.
\end{proof}

\begin{theorem}[Calibration by an integrated collision mean]
\label{thm:unk-scale-flux}
Suppose $2t+2R_*R_K<d_0$. The exact observed functions
$(g_1,g_2,F_1)$ determine the unknown ratio, the first physical
footprint in laboratory coordinates, and every protected original
component. In particular the two observed reciprocal mean pairs
suffice. Put $Q=-A$. For a matched pair $P_1=C+Q$, $P_2=C+rQ$, put
$p_D=p_{P_2}-p_{P_1}$ and use the flux $\Phi_C$ of
Lemma~\ref{lem:unk-scale-flux}. Then
\begin{equation}\label{eq:unk-scale-flux-inverse}
 a=\frac{\mathcal W(P_1)-2\Phi_C/t}{\mathcal W(D)},
 \qquad r=1+a^{-1},\qquad
 p_Q=a p_D,\qquad p_C=p_{P_1}-a p_D.
\end{equation}
The inverse is uniformly stable: expanded-support errors $e$ in
supremum norm and $u$ in $C^2$, together with flux error $m$, give
ratio and Hausdorff errors $O(e+m)$ and laboratory $C^2$ errors
$O(u+m)$. These assertions hold for sufficiently small $u+m$,
with constants depending only on the fixed priors.
\end{theorem}
\begin{proof}
The stopped inverse recovers the complete positivity supports from
$g_1,g_2$, and their fixed nesting margin gives their pairing.
Since $D=(r-1)Q$, one has $P_1=C+aD$ with $a=(r-1)^{-1}$.
Width additivity and \eqref{eq:unk-scale-flux} yield the first formula
in \eqref{eq:unk-scale-flux-inverse}; the remaining formulas follow
by support addition. The denominator is bounded below by $4g_*r_K$.
Also $a$ belongs to the fixed interval
$[1/(R_*-1),1/g_*]$. Consequently errors $e$ in the support values
and $m$ in the flux change $a$ and $r$ by at most $C(e+m)$.
The reconstructed supports then have supremum error $O(e+m)$ and
$C^2$ error $O(u+m)$. Their uniform positive curvature margins
persist at sufficiently small errors, so they represent physical
strictly convex bodies. Reflection recovers $A$ from $Q=-A$.
No first harmonic is discarded, and the original components retain
their laboratory locations.
\end{proof}

\begin{proposition}[A finite integrated collision measurement]
\label{prop:unk-scale-flux-measure}
Assume $2t+2R_*R_K<d_0$. After a prior-dependent coarse acquisition
of one complete first-setting component $P_1$, its integrated
forward mean $\Phi_C$ can be estimated to error $m$ with conditional
confidence $1-\varepsilon$ using
\begin{equation}\label{eq:unk-scale-flux-cost}
                       C m^{-2}\log(C/\varepsilon)
\end{equation}
attempted pooled collision bits. The same order bounds the number
of nominal centers. A dyadic nominal mesh of width $c m$ suffices,
uniformly over the possibly unbounded launch density.
\end{proposition}
\begin{proof}
Let $d_*>2t$ be the lower separation of the first-setting expanded
components. Choose a fixed $b<(d_*-2t)/8$. A sufficiently accurate
coarse hull determines a domain $E$, which is a finite union of
fixed dyadic squares, such that
\[
 P_1+t\overline B\subset E
       \subset P_1+(t+3b)\overline B.
\]
Choose all squares meeting an enlargement of the coarse hull;
reducing its fixed error and the fixed square size gives these
inclusions with positive margins. No other component's forward
response support meets $E$, since each such support is contained
in $P'_1+t\overline B$ and $d_*-2t-3b>0$.
Consequently $F_1=F_{1,C}$ on $E$ and
$\int_E F_1=\Phi_C$.

For $X$ uniform on $E$, one forward bit $Y$ at $X$ gives the bounded
variable $|E|Y$, whose expectation is $\Phi_C$. Independent fresh
repetitions and Hoeffding's inequality give the asserted count.
This uses the pooled forward bit itself and does not require a
direction-resolved mean.

To give the commands finite descriptions, subdivide the fixed
squares by a dyadic mesh $\zeta$ and sample their fine midpoints
uniformly. The expectation changes by $O(\zeta)$ without any
pointwise regularity assumption on $j$. Indeed, condition first
on the launch displacement and a compass direction. The hit
indicator is the difference between a convex segment sweep and
the original convex body. Intersect both with the fixed squares
of $E$. Their total boundary length and their number are uniformly
bounded. The midpoint sum differs from the integral only on fine
squares meeting these boundaries, whose total area is
$O(\zeta)$ by the convex boundary-tube estimate. The bound is
uniform in the launch displacement; averaging over its probability
law and over the four directions preserves it. Take
$\zeta\le c m$ and reserve half the error for this bias. The
confidence assertion is conditional on the coarse acquisition
and uses only subsequent fresh preparations.
\end{proof}

\subsection{Uniform finite reconstruction}
\begin{theorem}[Finite recovery without scale values]
\label{thm:unk-scale-finite}
Fix the numerical bounds of $\mathcal B_\eta$, the footprint and
boundary-mass bounds and \eqref{eq:unk-scale-prior}. Choose the fixed
dyadic compass step to satisfy $2t+2R_*R_K<d_0$. A finite adaptive
controller, supplied with the
two setting labels but neither their ratio nor a footprint or density
oracle, recovers the ratio, the physical first footprint, and a
smooth strictly convex periodic table. For all sufficiently small
$\nu>0$ and $0<\delta<1/4$, its ratio error, footprint laboratory
$C^2$ support error, and geometric loss are at most $C\nu$, and its
primitive discrete data are correct, with probability at least
$1-\delta$ uniformly over the admissible experiments.

Put
\[
                    q_\gamma=
                    \frac{(\gamma+3/2)s+1}{s-2}.
\]
The attempted-bit and nominal-center counts satisfy
\begin{align}
 N_\nu&\le
 C\nu^{-q_\gamma}\log(C/\nu)\log(C/(\nu\delta))
       +C\nu^{-2}\log(C/\delta),\label{eq:unk-scale-attempts}\\
 J_\nu&\le
 C\nu^{-1/(s-2)}\log(C/\nu)
       +C\nu^{-2}\log(C/\delta).\label{eq:unk-scale-centers}
\end{align}
Since $6<s\le7$ and $\gamma\ge0$, one has $q_\gamma>2$, so the
second term of \eqref{eq:unk-scale-attempts} is absorbed by the first.
Both physical spreads stay fixed as $\nu$ tends to zero. The finest
nominal mesh has width $c\nu^{s/(s-2)}$. The list of nominal centers,
setting labels and repetition counts has binary description length
at most
\begin{equation}\label{eq:unk-scale-control-description}
 C\left(\nu^{-1/(s-2)}\log(C/\nu)
              +\nu^{-2}\log(C/\delta)\right)
                         \log(C/(\nu\delta)).
\end{equation}
This is a bound on the digital control description. The fixed
homothety mechanism, its common origin, the known prior bounds,
and the physical realization of the nominal mesh remain apparatus
assumptions.
\end{theorem}
\begin{proof}
Run the rare-event boundary construction of
Proposition~\ref{prop:rare-query} and
Theorem~\ref{thm:rare-stationary} at both settings, stopping before
footprint subtraction. No step uses the scale value: the known
bounds on $r$ give the requisite component, curvature, separation,
and boundary-mass constants. With
$h\asymp\nu^{1/(s-2)}$ and $e=c h^s$, the two collections have
support errors $O(e)$ in supremum norm and $O(\nu)$ in $C^2$.
The fixed nesting margin pairs all complete components needed in
the protected patch. Split the confidence allowance between these
constructions and an independent integrated collision measurement.

Use one complete first-setting component and
Proposition~\ref{prop:unk-scale-flux-measure} with $m=c\nu$ to
estimate $\Phi_C$. Evaluate the four support values giving each
width sum, and apply \eqref{eq:unk-scale-flux-inverse}. The fixed
positive denominator and the stability statement of
Theorem~\ref{thm:unk-scale-flux} give $O(\nu)$ error in $r$ and
in the physical footprint. Subtract this one footprint from
every first-setting expanded support. The resulting supports have
$C^2$ and Hausdorff error $O(\nu)$, and retain strictly positive
curvature. The finite smoothing and period-locking constructions
apply after reducing $\nu_0$. Thus the output is a physical periodic
table with the required primitive data, real period coordinates,
and free area.

The two geometric reconstructions use
$C h^{-1}\log(C/h)$ local queries, each costing at most
$C e^{-(\gamma+3/2)}\log(C/(h\delta))$ attempts. The integrated
measurement adds $C\nu^{-2}\log(C/\delta)$ attempts and at most
that many centers. This proves both counts. Fresh preparations
give the required conditional estimates and a union bound gives
the stated confidence.

Choose the common fine dyadic mesh comparable to $e$ and dividing
the fixed coarse mesh in the integrated collision construction. Since $e\le c\nu$,
it also meets that construction's bias reserve. Every nominal point
lies in a fixed aperture and therefore needs $O(\log(C/\nu))$
digits. Every repetition count is polynomial in $1/\nu$ times a
logarithmic confidence factor and needs at most
$O(\log(C/(\nu\delta)))$ digits. Encoding the finite command list
therefore gives \eqref{eq:unk-scale-control-description}.
\end{proof}

\subsection{An independent area normalization}
For a component $C$, put $Q=-A$ and
\[
 P_1=C+Q,\qquad P_2=C+rQ,
 \qquad v_{1,C}(x)=\int_Aj(z)\one_C(x+z)\dd z.
\]
The subscript $C$ denotes a single complete positive component. The
stopped inverse recovers this restriction of the occupation without
knowing $j$. Fubini's theorem gives the additional invariant
\begin{equation}\label{eq:unk-scale-mass}
 M_C:=\int_{P_1}v_{1,C}(x)\dd x
     =\int_Aj(z)|C-z|\dd z=|C|.
\end{equation}
In particular, the area on the right is measured by the collision
experiment; it is not an input concerning the obstacle.

For planar convex bodies $E,F$, let $V(E,F)$ be their mixed area,
with $V(E,E)=|E|$. In support coordinates,
\begin{equation}\label{eq:unk-scale-mixed-area}
 |E|=\frac12\int_0^{2\pi}(p_E^2-(p_E')^2)\dd\theta,
 \qquad
 V(E,F)=\frac12\int_0^{2\pi}(p_Ep_F-p_E'p_F')\dd\theta.
\end{equation}
These expressions are unchanged by translation. Support addition and
the quadratic area formula are the usual mixed-area identities
\cite{Schneider}.

\begin{theorem}[An unknown scale ratio is identifiable]
\label{thm:unk-scale-exact}
Under the initial condition $t+2R_*R_K<d_0$, the exact reciprocal
differences $(g_1,g_2)$ alone determine $r$, the physical footprint
$A$ in laboratory coordinates,
and every protected original component. The full period group is then
determined by the retained exact recognition theorem.

More explicitly, match a complete pair $P_1,P_2$ by nesting and let
$D$ be the convex body with support
\begin{equation}\label{eq:unk-scale-difference}
                         p_D=p_{P_2}-p_{P_1}.
\end{equation}
Set
\[
 a_0=|P_1|,\qquad b_0^*=V(P_1,D),\qquad c_0=|D|,
 \qquad \mathcal D=(b_0^*)^2-c_0(a_0-M_C).
\]
Then $c_0>0$, $\mathcal D>0$, and
\begin{align}
 a&=\frac{b_0^*-\sqrt{\mathcal D}}{c_0}
    =\frac{a_0-M_C}{b_0^*+\sqrt{\mathcal D}},
 &r&=1+a^{-1},\label{eq:unk-scale-root}\\
 p_Q&=a p_D,
 &p_C&=p_{P_1}-a p_D.\label{eq:unk-scale-inverse}
\end{align}
All support functions in these formulas are in the laboratory frame.
No centering convention for $A$ or for its density is imposed.
\end{theorem}
\begin{proof}
Lemma~\ref{lem:scale-matching} uses only the lower nesting margin and
the separation, so it supplies the pairing without the value of $r$.
For the true pair, $D=(r-1)Q$ and $a=(r-1)^{-1}$. In particular $D$
is a full-dimensional strictly convex body containing a fixed inball.
The equation $P_1=C+aD$ and the mixed-area identity give
\[
           M_C=a_0-2a b_0^*+a^2c_0.
\]
Moreover,
\begin{equation}\label{eq:unk-scale-root-margin}
 b_0^*-ac_0=V(C,D)>0,
 \qquad \mathcal D=V(C,D)^2.
\end{equation}
Thus the true $a$ is precisely the smaller root in
\eqref{eq:unk-scale-root}. Any other positive value $a'$ for which
$p_{P_1}-a'p_D$ is the support of a convex body of area $M_C$ would
also satisfy $b_0^*-a'c_0>0$, and hence would be the same smaller
root. This proves uniqueness among physical candidates.
Formula~\eqref{eq:unk-scale-inverse} now recovers $C$ and $Q$;
reflection recovers $A$. The answer for $r$ and $A$ is common to all
matched pairs. The recognition theorem applies to the recovered
original components.
\end{proof}

\begin{proposition}[Stability of the area-normalized inverse]
\label{prop:unk-scale-stability}
Under the fixed priors above, suppose the matched expanded supports
at both settings are known with supremum error $e$ and $C^2$ error $u$, and
$M_C$ is known with error $m$. For sufficiently small $u+m$, the
smaller-root construction is well defined and gives
\begin{align}
 |\widehat r-r|+d_H(\widehat A,A)
                 +\max_Cd_H(\widehat C,C)&\le C(e+m),
                                      \label{eq:unk-scale-stability-0}\\
 \|p_{\widehat A}-p_A\|_{C^2}
       +\max_C\|p_{\widehat C}-p_C\|_{C^2}&\le C(u+m).
                                      \label{eq:unk-scale-stability-2}
\end{align}
The constants depend on the stated bounds, including $g_*$, but not
on the unknown ratio, footprint, or density.
\end{proposition}
\begin{proof}
The true $D$ has uniformly positive curvature radius. Hence the
estimated support difference remains a strictly convex support for
small $u$. On bounded convex classes area and mixed area are Lipschitz
in Hausdorff distance, so their estimated coefficients have error
$O(e)$. For example, area follows from the two parallel-body
inclusions, and mixed area follows by polarization using $|E+F|$.

The inball of $D$ and the lower perimeter bound for $C$ imply
$V(C,D)\ge c>0$. Equation~\eqref{eq:unk-scale-root-margin} therefore
keeps the discriminant uniformly away from zero. The two equivalent
root formulas have bounded derivatives throughout a fixed
neighborhood of the admissible data. Also
$1/(R_*-1)\le a\le1/g_*$. Consequently
$|\widehat a-a|+|\widehat r-r|\le C(e+m)$.
Multiplication by $p_D$, whose $C^2$ norm is uniformly bounded in the
protected frame, proves both estimates. The positive curvature
margins of the original bodies and $A$ persist, so the recovered
supports represent actual convex bodies. A single pair may be used
for $a$ and $A$, after which the same footprint is subtracted from all
other first-setting components.
\end{proof}

\subsection{A finite measurement of one component area}
The area normalization can be measured without a fine deterministic
grid of occupation queries. The following construction uses a fixed
coarse domain and a bounded weight computable before the area
measurements begin.

\begin{proposition}[Integrated occupation from pooled bits]
\label{prop:unk-scale-mass}
After a prior-dependent coarse acquisition of one complete expanded
component $P_1$, its mass $M_C$ can be estimated to error $m$ with
conditional confidence $1-\varepsilon$ using
\begin{equation}\label{eq:unk-scale-mass-cost}
                       C m^{-2}\log(C/\varepsilon)
\end{equation}
attempted pooled collision bits at the first setting. The same order
bounds the number of nominal centers. A dyadic nominal mesh of width
$c m$ suffices. All centers lie in one fixed protected aperture, and
the assertion is uniform over the possibly unbounded density $j$.
\end{proposition}
\begin{proof}
Let $d_*>t$ be the known lower separation of the expanded components.
Choose a fixed positive $b<(d_*-t)/8$. From a sufficiently accurate
coarse hull construct a domain $E$, which is a finite union of squares
of a fixed dyadic mesh $d$, such that
\begin{equation}\label{eq:unk-scale-domain}
 P_1\subset E,\qquad
 \dist(P_1,E^c)>b,\qquad E\subset P_1+3b\overline B.
\end{equation}
The mesh may be chosen to divide the dyadic compass step $t$.
Taking all mesh squares meeting a fixed small enlargement of the
coarse hull gives these inclusions, after reducing its fixed error
and $d$. In particular no other positive component meets
$E+t\overline B$. Boundaries of squares may be assigned consistently
to one neighboring square; they have measure zero.

On $E$ the first-setting forcing therefore satisfies
\[
             g_1=(T^E-I)v_{1,C}.
\]
Let $w_E=(I-T^E)^{-1}\one_E$. The killed-walk exit bound gives
$0\le w_E\le H_E$, where $H_E$ is a fixed prior constant. Because
$d$ divides $t$, the state graph of the killed walk is identical for
all starting points in the interior of any one mesh square. Thus
$w_E$ is constant on each of the finitely many squares. Its values
are the solution of a finite rational linear system with coefficients
$0,1,1/4$, and are available without any occupation observations.
The inverse exists since the walk exits the bounded domain almost
surely and has uniformly bounded mean exit time.

The compass kernel is symmetric. Integrating the killed operator on
$E$ and interchanging its finitely many translations gives
\begin{equation}\label{eq:unk-scale-duality}
 \int_E w_E g_1
   =\int_E v_{1,C}(T^E-I)w_E
   =-\int_Ev_{1,C}=-M_C.
\end{equation}
If $X$ is uniform on $E$, take one fresh forward bit $Y^+$ and one
fresh reciprocal reverse bit $Y^-$ at $X$. The variable
\[
            Z_E=-|E|w_E(X)(Y^+-Y^-)
\]
has expectation $M_C$ and absolute value at most $|E|H_E$. Independent
repetitions and Hoeffding's inequality prove
\eqref{eq:unk-scale-mass-cost} for continuous nominal sampling.
This calculation does not assume that an occupation bit is observed.

For a finite command description choose a dyadic mesh $\zeta$ dividing
$d$ and sample $X_\zeta$ uniformly from the fine square midpoints in
$E$. This changes the expectation by $O(\zeta)$, uniformly over $j$.
Here is a direct bound which does not require an upper bound or a
modulus of continuity for the density. Before averaging in the
launch displacement, each forward or reverse bit is a finite
difference of indicators of convex obstacles and their segment
sweeps, translated by the displacement. Intersect these sets with
the fixed squares of $E$, on which $w_E$ is constant. Their total
boundary length and number are uniformly bounded. A midpoint sum
and its integral agree on every fine square which does not meet
one of these boundaries. The total area of the remaining squares is
$O(\zeta)$, by the boundary-tube estimate for convex sets. The bound
is uniform in the launch displacement, so integration against the
probability density preserves it. The same argument applies to the
four pooled directions and the reciprocal shifts. Taking
$\zeta\le c m$ reserves at most half the permitted error for this
quadrature bias. Hoeffding's inequality supplies the other half.
The procedure and its confidence bound hold conditional on the
successful coarse acquisition, using fresh independent preparations.
\end{proof}

The normalization argument uses the first setting's density only
through its integral and the quantitative boundary lower bound.
It therefore remains valid if the two settings have different
stationary densities on $A$ and $rA$, provided their normalized
boundary-mass bounds are uniform. Their supports must still be
exactly homothetic. Finally, an unknown common unscaled displacement
$b$ in the commands replaces $C$ by $C-b$ throughout, leaving $r$
and the physical footprint unchanged. Translating the table and
$b$ by the same vector leaves every collision bit unchanged. The
absolute location in that extended experiment requires a laboratory
reference.
